\pdfoutput=1
\documentclass[11pt]{article}

\usepackage[margin=1.05in]{geometry}
\usepackage[utf8]{inputenc}
% Latin Modern with T1 when available (arXiv has it); plain CM otherwise
\IfFileExists{lmodern.sty}{\usepackage[T1]{fontenc}\usepackage{lmodern}}{}
\usepackage{amsmath,amssymb,amsthm}
\usepackage{booktabs}
\usepackage{microtype}
\usepackage{enumitem}
\usepackage[hidelinks]{hyperref}

\newtheorem{theorem}{Theorem}
\newtheorem{corollary}{Corollary}
\newtheorem{proposition}{Proposition}
\newtheorem{lemma}{Lemma}
\theoremstyle{definition}
\newtheorem{definitionx}{Definition}

% Shared macro definitions for paper.tex and paper.fr.tex.
% One source of truth: a symbol renamed here is renamed in both papers.
\newcommand{\Fp}{\mathbb{F}_p}
\newcommand{\negl}{\mathsf{negl}}
\newcommand{\LeafHash}{\mathsf{LeafHash}}
\newcommand{\Compress}{\mathsf{Compress}}
\newcommand{\PRF}{\mathsf{PRF}}
\newcommand{\CtxHash}{\mathsf{CtxHash}}
\newcommand{\Trunc}{\mathsf{Trunc}_4}
\newcommand{\MerkleFold}{\mathsf{MerkleFold}}
\newcommand{\RO}{\mathsf{RO}}
\newcommand{\Enc}{\mathsf{Enc}}
% Poseidon2 state width is t, so the Fiat-Shamir challenge vector is
% chi; the per-threshold accumulators are calligraphic, so they do not
% collide with the transcript body T.
\newcommand{\chall}{\chi}
\newcommand{\tree}{\mathcal{T}}
\newcommand{\adv}{\mathcal{A}}
% := without pulling in mathtools; arXiv has it, but one fewer package
% is one fewer thing that can differ between their TeX Live and yours.
\providecommand{\coloneqq}{\mathrel{\vcenter{\hbox{$:$}}}\!=}


\title{Post-Quantum Anonymous Set Membership for\\
Surveillance-Free Online Age Verification}
\author{Alexandre Poumaroux\thanks{\texttt{research@kidev.org}}}
\date{}

\begin{document}
\maketitle

\begin{abstract}
We present a protocol that lets websites verify that a visitor is above the age of
majority while making it cryptographically impossible for either the verifying website or
the issuing government to learn who the visitor is or which sites any citizen visits. The
core construction is anonymous set membership: the government publishes a single signed
Merkle root committing to the set of hash commitments of verified adults, and each
citizen's open-source wallet proves in zero knowledge that it knows the preimage of some
leaf under that root, together with the correct derivation of a per-site, per-epoch
pseudonymous key. Verification is performed entirely offline by the website against the
published root; the government is structurally absent from every usage flow, so there
exists no verification traffic, timestamp, or counter that could be logged, correlated, or
subpoenaed. All primitives are hash based or lattice based: Poseidon2 over the Goldilocks
field for commitments, Merkle compression, and pseudonym derivation; the MPC-in-the-head
proof system ZKBoo compiled with Fiat-Shamir for the zero-knowledge argument; and
ML-DSA-65 (FIPS 204) for root authentication. The design consequently withstands quantum
adversaries. We give game-based definitions and complete proofs for anonymity,
unforgeability, and replay resistance, an honest treatment of the one property no
credential system can enforce (voluntary lending), and measured performance of a complete
reference implementation: at 128-bit soundness over a tree with 4.29 billion leaf
capacity, proving takes 4.7\,s and verification 3.0\,s on a single core of an AMD Ryzen~9 7950X in pure
Python, with a 21.8\,MB proof, while the accompanying Rust port of the same backend proves
in 0.33\,s and verifies in 0.18\,s on one core, fifteen times faster than CPython measured
on that same machine. We analyze the path to sub-200\,KB proofs and tens-of-milliseconds
verification with production proof backends.
\end{abstract}

\section{Introduction}

Online age verification mandates are proliferating, for social networks, app stores, and
a growing range of age-gated services, and every widely deployed compliance mechanism
forces a fatal trade: either the website learns who you are (document upload, face scan),
or a third party learns where you go (centralized verification services, federated
identity). Both create databases whose mere existence is the harm: a standing record of
who visits what online is general-purpose surveillance infrastructure, one leak, one
subpoena, or one policy change away from being turned against any citizen or group. The
result is predictable: adults evade the checks, minors follow the same evasion routes, and
the stated child-protection goal fails while the surveillance risk remains.

We take as our design constraint the strongest possible version of the privacy
requirement, stated as an inviolable rule: the system must let a website confirm ``this
visitor is an adult certified by the government'' while giving neither the website nor the
government, nor the two colluding, any capability, ever, of mapping a visit to a citizen.
The guarantee must be mathematical, not institutional: a citizen who assumes the
government is actively malicious must still be safe using the system.

Formally, the system must satisfy:
\begin{description}[leftmargin=2.6em,style=nextline]
\item[R1 (Soundness).] A website accepts only visitors holding a credential issued by the
government to a verified adult.
\item[R2 (Issuer blindness).] The government cannot determine which citizen is behind any
verification event, nor that any particular citizen used the system at all after
enrollment.
\item[R3 (Verifier blindness).] The website learns exactly one bit (adult: yes) plus,
optionally, a pseudonym stable only within that site and a bounded time window; pseudonyms
across sites are unlinkable even to colluding sites.
\item[R4 (Collusion resistance).] R2 and R3 hold even if the government and any set of
websites pool all their data.
\item[R5 (Post-quantum security).] All properties hold against quantum adversaries.
\end{description}

Two structural decisions do most of the work. First, the government signs no per-citizen
value at any point: its entire cryptographic output is one signature on one public
accumulator, so no signature can ever act as a tracking tag. Second, verification requires
only public parameters, so it is performed locally by the website; the government has no
verification endpoint, and messages that are never sent cannot be logged, correlated, or
subpoenaed.

Our contributions:
\begin{enumerate}[leftmargin=2em]
\item A protocol built from anonymous set membership that satisfies R1 through R5, in
which the government's entire online footprint is one signed public value and verification
is a local computation at the website (Section~\ref{sec:protocol}).
\item Game-based security definitions and complete proofs (Section~\ref{sec:security}),
including an exact accounting of what each party and each colluding coalition can compute,
and an explicit statement of the properties the system does not provide.
\item A complete reference implementation at production parameters:
Poseidon2-Goldilocks with the official constants of its designers, 256-bit digests, a
$2^{32}$-capacity sparse Merkle accumulator, $2^{-128}$ Fiat-Shamir soundness, ML-DSA-65
root signing, an encrypted wallet format, and a binary wire format; a native Rust port of
the prover and verifier interoperable with it at the wire level; plus measured benchmarks
and a test suite covering the attacks analyzed in Section~\ref{sec:security}
(Sections~\ref{sec:impl}, \ref{sec:perf}).
\end{enumerate}

\section{Related Work}

\paragraph{Anonymous credentials.} Chaum's blind signatures \cite{chaum} began the study
of credentials that issuers cannot trace. Brands \cite{brands} and Camenisch-Lysyanskaya
\cite{cl} built full anonymous credential systems; modern deployments (Idemix, BBS+ in
verifiable-credential profiles \cite{w3cvc}) descend from the latter. All practical
instantiations rely on RSA, discrete-log, or pairing assumptions and fall to Shor's
algorithm \cite{shor}; post-quantum anonymous credentials and blind signatures remain
research prototypes with large keys, large signatures, restricted concurrent security, or
non-standard assumptions \cite{aksy,delpino}.

\paragraph{Set-membership pseudonymity.} Architecturally, our construction descends
from Semaphore-style anonymous signaling \cite{semaphore} and from the nullifier design of
Zerocash \cite{zerocash}: identity is a secret whose commitment lives in a public Merkle
set, membership is proven in zero knowledge, and deterministic PRF outputs provide scoped
pseudonyms (external nullifiers). We adapt this pattern to age verification, replace its
discrete-log-based proof systems and hashes with post-quantum components, and add
verifier-side transcript binding for replay resistance without any coordinating server.

\paragraph{Rate-limited anonymous tokens.} Privacy Pass \cite{privacypass} and its
rate-limited variants let users spend unlinkable tokens; they address a different problem
(per-request unlinkability with an online issuer) and their VOPRF foundations are
pre-quantum.

\paragraph{Age verification systems.} Commercial providers (document upload, face
estimation) reveal identity or biometrics to a party that must be trusted. The W3C
Verifiable Credentials model \cite{w3cvc} supports selective disclosure but its
unlinkable-presentation profiles rest on BBS+ (pairings, hence pre-quantum), and its
issuer-involved presentation flows leak usage metadata. The EU Digital Identity Wallet's
age-attestation work shares our goal for R3 but not our threat model for R2/R4: it trusts
the issuer not to correlate.

\paragraph{Proof systems.} Zero-knowledge proofs originate with Goldwasser, Micali, and
Rackoff \cite{gmr}. ZKBoo \cite{zkboo} and its successors (ZKB++/Picnic \cite{picnic}, KKW
\cite{kkw}) construct zero-knowledge arguments from the MPC-in-the-head paradigm of Ishai
et al.\ \cite{ikos}, using only symmetric primitives: no trusted setup, plausible
post-quantum security. STARKs \cite{stark} achieve the same trust profile with succinct
proofs at the cost of implementation complexity. Our reference implementation uses ZKBoo
for auditability; Section~\ref{sec:deploypath} quantifies the STARK/VOLE-in-the-head
deployment path.

\paragraph{Post-quantum standards.} NIST standardized ML-DSA (Dilithium) as FIPS 204
\cite{fips204,dilithium} and SHA-3 as FIPS 202 \cite{fips202}; Poseidon2 \cite{poseidon2}
is the current standard arithmetization-friendly permutation for zero-knowledge systems.
Grover's algorithm \cite{grover} halves symmetric security levels, which our parameter
choices account for.

\section{System and Threat Model}
\label{sec:model}

\paragraph{Participants.} A Government Authority (GA) that verifies age through existing
civil registry mechanisms; Citizens running an open-source Wallet on their own devices;
Websites that must gate content on adulthood.

\paragraph{Trust assumptions.} (i) The GA verifies age correctly at enrollment and enrolls
only adults; (ii) cryptographic primitives satisfy their standard security notions
(collision resistance and PRF security of the Poseidon2-based functions of
Section~\ref{sec:prims}, EUF-CMA of ML-DSA-65, random-oracle behavior of SHA3/SHAKE
\cite{bellrog} for commitments and Fiat-Shamir); (iii) the citizen's device is not
compromised. We explicitly do \emph{not} assume the GA refrains from correlating,
colluding, or coercing; that websites behave; or that citizens keep credentials to
themselves.

\paragraph{Adversaries.} A malicious GA may record everything it ever sees, collude with
websites, and attempt split-view attacks on published data. Malicious websites may pool
data with each other and the GA and attempt cross-site linkage. Malicious citizens may
attempt to forge membership, replay tokens, or lend credentials. All adversaries may be
quantum-capable.

\paragraph{Security goals.}
\textbf{G1} Unforgeability: no coalition lacking an enrolled secret produces an accepted
token, except with negligible probability.
\textbf{G2} Anonymity: for any accepted token, no coalition of GA and websites can
distinguish which of the enrolled citizens produced it with advantage non-negligibly
better than guessing among all enrolled citizens consistent with public data.
\textbf{G3} Cross-site unlinkability: tokens presented at different sites (or in different
epochs) cannot be linked to the same citizen.
\textbf{G4} Replay resistance: a token accepted once is not accepted again, at the same or
any other site.
\textbf{G5} All of the above against quantum adversaries.

\paragraph{Out of scope.} Network-layer identifiers (IP addresses) are visible to websites
as with any HTTPS request; citizens who require network anonymity compose the system with
Tor or a VPN, which it does not interfere with, precisely because no government endpoint
is contacted. Device compromise and legal compulsion of the citizen herself are out of
scope.

\section{The Protocol}
\label{sec:protocol}

\subsection{Primitives and parameters}
\label{sec:prims}

\paragraph{Field.} The Goldilocks prime $p = 2^{64} - 2^{32} + 1$. All protocol values are
vectors over $\Fp$; digests are 4 elements (256 bits).

\paragraph{Permutation.} Poseidon2 \cite{poseidon2} with state width $t = 8$, S-box $x^7$,
8 external and 22 internal rounds, using the round constants and internal matrix published
by its designers. Throughout the paper, $P : \Fp^8 \to \Fp^8$ denotes this permutation and
$\Trunc : \Fp^8 \to \Fp^4$ the projection of a state onto its first four elements. From
$P$ we build three domain-separated functions, using three \emph{distinct} field
constants $\mathit{tag}_{\mathrm{leaf}}, \mathit{tag}_{\mathrm{prf}},
\mathit{tag}_{\mathrm{empty}} \in \Fp$ (their pairwise distinctness is what makes the
domain separation load-bearing in Theorem~\ref{thm:unforge}):
\begin{align*}
\LeafHash(s) &= \Trunc\bigl(P(s_0, s_1, s_2, s_3, \mathit{tag}_{\mathrm{leaf}}, 0, 0, 0)\bigr),\\
\Compress(l, r) &= \Trunc\bigl(P(l \,\|\, r) + (l \,\|\, r)\bigr),\\
\PRF(s, c) &= \Trunc\bigl(P(P(s_0, s_1, s_2, s_3, \mathit{tag}_{\mathrm{prf}}, 0, 0, 0) \boxplus c)\bigr),
\end{align*}
the enrollment commitment of secret $s \in \Fp^4$, the 2-to-1 Merkle compression (the
feed-forward compression mode recommended by the Poseidon2 authors), and the pseudonym
derivation. Here $l \,\|\, r \in \Fp^8$ denotes concatenation and $x \boxplus c \coloneqq
(x_0 + c_0, \ldots, x_3 + c_3, x_4, \ldots, x_7)$ adds the 4-element context $c$ into the
rate, leaving the capacity untouched. Contexts are $c
= \CtxHash(\mathrm{domain}, \mathrm{epoch})$, 4 elements derived by SHAKE256 with
rejection sampling; multi-threshold deployments fold the age threshold into the hashed
domain string (Section~\ref{sec:thresholds}). Under standard sponge and compression-mode analyses these provide
128-bit collision resistance and 128-bit post-quantum preimage resistance (capacity and
digest are 256 bits; Grover \cite{grover} halves the exponent).

\paragraph{Accumulator.} A sparse incremental Merkle tree \cite{merkle} of depth $d = 32$
(capacity $2^{32} \approx 4.29$ billion, sufficient for any national or federated
population) with all-equal empty leaves $E =
\Trunc(P(0,0,0,0,\mathit{tag}_{\mathrm{empty}},0,0,0))$ and cached empty-subtree digests,
so insertion and path computation cost $O(d)$ regardless of occupancy.

\paragraph{Zero-knowledge argument.} ZKBoo \cite{zkboo}, an MPC-in-the-head
\cite{ikos} argument system realizing zero-knowledge proofs \cite{gmr}, over the
arithmetic circuit of Section~\ref{sec:issuance}, with the $(2,3)$ additive decomposition
over $\Fp$, party tapes expanded from 256-bit seeds by SHAKE256 in counter mode, SHA3-256
view commitments, and Fiat-Shamir challenges derived by SHA3-256 over the full transcript,
modeled as a random oracle \cite{bellrog}. Each repetition has soundness error $2/3$; we
use $\tau = 219$ repetitions, giving error $(2/3)^{219} < 2^{-128}$. Input shares of two
parties are tape-derived (the ZKB++ optimization \cite{picnic}); the third party's shares
and the neighbor view's multiplication outputs are the only per-gate data transmitted.

\paragraph{Root authentication.} ML-DSA-65 (FIPS 204 \cite{fips204}), NIST security
category 3, signs each published root bundle.

\paragraph{Wallet storage.} The secret is encrypted at rest with AES-256-GCM under a key
derived by scrypt ($N = 2^{15}$, $r = 8$, $p = 1$).

\subsection{Enrollment}

The wallet samples $s \leftarrow \Fp^4$ uniformly (256 bits of local entropy that never
leave the device) and computes $\mathit{leaf} = \LeafHash(s)$. The citizen authenticates
to the GA by whatever existing civil-identity mechanism the jurisdiction uses (in person
or via eID) and submits only $\mathit{leaf}$. The GA verifies the citizen is an adult not
previously enrolled, appends $\mathit{leaf}$ to the tree at the next free index, and
records $(\mathrm{citizen\_id}, \mathrm{index})$. Key generation is entirely local:
producing $s$ and $\mathit{leaf}$ requires no network connection at all, so a wallet app
can create its credential fully offline (even on an air-gapped device) and disclose the
commitment only at the moment of enrollment; nothing about the key ever depends on a
server.

The GA is permitted to remember this mapping forever. It is harmless: $\mathit{leaf}$ and
index never appear in any later message the GA could observe. Anonymity does not rest on
the GA forgetting; it rests on the GA never seeing anything else.

\subsection{Root publication}
\label{sec:rootpub}

After each batch of enrollments or revocations the GA publishes a root bundle
\[
B = \bigl(\mathit{root}, \mathit{size}, \mathit{issued}, d,\
\mathrm{Sig}_{\text{ML-DSA-65}}(\ell \,\|\, \mathit{root} \,\|\, \mathit{size} \,\|\,
\mathit{issued} \,\|\, d)\bigr),
\]
where $\mathit{size}$ is the number of leaves allocated so far, $\mathit{issued}$ the
publication timestamp, and $\ell$ is a fixed version label (extended with the age threshold in the
multi-threshold deployment of Section~\ref{sec:thresholds}), published together with the
full leaf sequence, the tree being public data. Bundles are distributed
transparency-log style: mirrored, gossiped, and append-only audited, so the GA cannot
serve different roots to different users (a split-view attack would otherwise let it give
one citizen a personalized tree and recognize the resulting proofs; gossip reduces this to
an attack detectable by comparing two bundle observations). Wallets synchronize the tree
in bulk, or via private information retrieval, so obtaining one's path reveals nothing
about which leaf is one's own; websites need only the bundle.

\subsection{Token issuance}
\label{sec:issuance}

To visit domain $D$ in epoch $e = \lfloor \mathrm{now} / 30\ \mathrm{days} \rfloor$, the
wallet obtains a fresh 128-bit nonce $n$ from the website, computes $c = \CtxHash(D, e)$,
the pseudonym $k = \PRF(s, c)$, the Merkle path $(b_0, \sigma_0, \ldots, b_{d-1},
\sigma_{d-1})$ for its leaf, where $b_i \in \{0,1\}$ is the $i$-th bit of the leaf index
and $\sigma_i \in \Fp^4$ the sibling digest at level $i$, and a proof $\pi$ of the
statement
\begin{align*}
S(\mathit{root}, k, c):\quad \exists\, s, (b_i, \sigma_i)_{i<d}:\quad
& b_i \in \{0,1\},\\
& \MerkleFold\bigl(\LeafHash(s), (b_i, \sigma_i)_{i<d}\bigr) = \mathit{root},\quad
\PRF(s, c) = k.
\end{align*}
Here $\MerkleFold$ is the iterated root recomputation defined by $h_0 = x$ and, for $0 \le
i < d$, $h_{i+1} = \Compress(h_i, \sigma_i)$ if $b_i = 0$ and $h_{i+1} =
\Compress(\sigma_i, h_i)$ if $b_i = 1$, with $\MerkleFold\bigl(x, (b_i,
\sigma_i)_{i<d}\bigr) = h_d$. The statement is realized as the arithmetic circuit: one
$\LeafHash$ permutation; per level, a 1-multiplication conditional swap per digest
element, one booleanity constraint $b_i^2 - b_i = 0$ exposed as a public zero output, and
one $\Compress$ permutation; and two permutations for the PRF. Circuit size is $344(d+3) +
5d$ multiplication gates (12{,}200 at $d = 32$). The Fiat-Shamir transcript hashes, in
addition to all commitments and output shares, the context string
\[
\mathit{ctx} = \text{``avsm-token/1''} \,\|\, D \,\|\, e \,\|\, \mathit{root} \,\|\, k
\,\|\, n,
\]
binding the proof to the domain, epoch, root, claimed pseudonym, and the website's nonce,
where consecutive fields are joined by a fixed reserved separator byte (the byte
\texttt{|}, which occurs in no hostname). Write $\Enc(D, e, \mathit{root}, k, n)$ for this
byte string; the replay argument of Section~\ref{sec:security} needs one syntactic property
of it.

\begin{lemma}[Context encoding]
\label{lem:enc}
If the domain field contains no occurrence of the separator byte, then $\Enc$ is
injective.
\end{lemma}

\begin{proof}
Every field after the domain has a fixed width: 8 bytes for $e$, 32 each for
$\mathit{root}$ and $k$, and $n$ runs to the end of the string. A separator-free domain
field is therefore delimited by the first separator, and the remaining fields are read off
at fixed offsets, so the byte string determines the tuple.
\end{proof}

\noindent
Hostnames contain no separator byte, but that is a property of the caller's input rather
than of the protocol, so wallet and verifier both reject a domain containing one before
hashing anything. The token is $(k, e, D, \pi)$.

\subsection{Verification (offline)}
\label{sec:verification}

The website checks that $n$ is one it issued, unexpired (600\,s TTL) and unconsumed; that
$D$ matches itself and $|e - e_{\mathrm{now}}| \le 1$, where $e_{\mathrm{now}}$ is the
epoch of the verifier's own clock; recomputes $c$ and the expected public outputs $(k,
\mathit{root}, 0^d)$, where $0^d$ is the vector of $d$ zeros the booleanity constraints
must produce using the root from its latest verified bundle;
and runs the proof verifier: for each repetition, derive the challenge from the
transcript, re-expand the two opened seeds, recompute the first opened party's entire view
and the second opened party's wire evaluation from its committed multiplication outputs,
and check both commitments, both parties' output shares, and the three-share
reconstruction of every public output. On success it marks $n$ consumed and may associate
$k$ with a session or account. No message leaves the website.

\subsection{Revocation and epochs}

Revocation (stolen wallet, erroneous enrollment, death) is a leaf reset to $E$ and a new
bundle; the revoked secret no longer satisfies $S$ against the new root, and all other
citizens' proofs continue to verify after their wallet syncs the updated path. Epochs
bound pseudonym lifetime: a site's key $k$ rotates every 30 days, so long-horizon profile
building at a single site is capped, while within an epoch the stable $k$ gives sites a
returning-visitor handle without cookies or accounts. Websites accept tokens whose epoch
is within one of their current epoch (the previous or the next, matching the check $|e -
e_{\mathrm{now}}| \le 1$ of Section~\ref{sec:verification}), tolerating clock skew and
boundary races around epoch rollover.

\subsection{Multiple age thresholds}
\label{sec:thresholds}

Nothing in the construction is specific to the age of majority; adulthood is simply the
predicate the GA checks at enrollment. A jurisdiction whose laws gate different services
at different ages (13, 15, 16, 18) runs one accumulator per threshold: the GA maintains
parallel trees $\tree_{13}, \tree_{15}, \ldots$, one per threshold $a$, each with its own root
$\mathit{root}_a$, inserts a citizen's commitment into every tree whose threshold the
citizen satisfies, and publishes one signed bundle per tree with $a$ bound under the
signature, so a root cannot be presented as certifying a different threshold than the one
it was published for. The same wallet secret and the same commitment serve every tree,
though the trees fill independently, so the wallet stores one leaf index per tree;
insertions for citizens crossing a threshold are batched with the regular bundle cadence;
revocation clears the leaf in every tree. All security theorems apply per tree unchanged,
and the anonymity set of a proof against $\tree_a$ is the set of enrolled citizens satisfying
threshold $a$, a population-scale set at every $a$.

\paragraph{The bracketing attack.} A website whose duty is threshold $a$ learns one bit,
age $\ge a$. Nothing forces a hostile website to stop there. A site that admits visitors
at 13 could additionally test each visitor against the 18 threshold: adults pass, a
fourteen-year-old fails, and the site has sorted the very users it admits into age
brackets. A multi-threshold design must close this channel, and must close it with
mathematics rather than with terms of service. Three mechanisms do so. The first two are
cryptographic; the third makes the only request they cannot eliminate visible to the
person it targets.

\begin{enumerate}[leftmargin=2em]
\item \emph{Threshold binding.} The threshold is welded into the proof itself. The
context becomes $c_a = \CtxHash(D, e, a)$ (implemented by injecting $a$ into the hashed
context string of Section~\ref{sec:issuance}), the pseudonym becomes $k = \PRF(s, c_a)$,
and the Fiat-Shamir transcript binds $(D, e, \mathit{root}_a, k, n)$ together with $a$. A
proof is therefore a proof \emph{for one threshold}: retesting it against any other
threshold's root fails, for every visitor, whether or not that visitor belongs to the
other tree (Proposition~\ref{prop:isolation}, Section~\ref{sec:isolation}). A retest that
fails for everyone reveals nothing about anyone.
\item \emph{Single use.} The transcript binds the site's fresh nonce and the verifier
consumes it (Theorem~\ref{thm:replay}), so a received token cannot be replayed against
anything, including another root. Probing therefore cannot recycle material the site
already holds; it would require asking the wallet for a brand new proof.
\item \emph{Wallet-side transparency.} That remaining request is made to the wallet, and
the wallet is the citizen's agent, not the site's. The wallet displays which threshold a
site asks it to prove, remembers the first threshold each domain declared, and treats any
change, or any request to prove a second threshold, as suspect: it interrupts with an
explicit warning, refuses by default, and proceeds only on deliberate user confirmation
(deployments may confine the override to advanced settings). The warning logic runs
before any membership-dependent computation, so it is triggered identically by every
wallet, whatever the holder's age; and a wallet holding no leaf in the requested
accumulator declines locally rather than emitting a proof that would fail, so the site
cannot distinguish that case from a refusal. The site cannot read an age off the wallet's reaction,
only off an answer the user knowingly chooses to give, and the probing request itself
becomes user-visible, recordable evidence of abuse.
\end{enumerate}

A domain's threshold thereby becomes part of its public identity. A site remains free to
declare 18 for a nominally 13+ service, but it then turns minors away uniformly and
visibly, and still receives one bit. The natural-seeming cascade (attempt $\tree_{18}$, fall
back on $\tree_{16}$, conclude $16 \le \mathrm{age} < 18$) is pointless for honest sites,
since a visitor proves directly against the site's one declared threshold, and is exactly
the double-request pattern that mechanism 3 flags. Because pseudonyms derive from the
thresholded context, tokens for different thresholds at one domain are additionally
unlinkable (Theorem~\ref{thm:unlink} applies, since $c_a \ne c_b$ for $a \ne b$), defense
in depth if a user ever does confirm a second threshold. The legal prohibition of
Section~\ref{sec:deploy}, point (iii), remains as a liability backstop, but the mechanism
does not depend on it. The reference implementation enforces all of this: bundles carry
the threshold under the GA signature, the context hash folds the threshold in, wallets
maintain a persistent pin store and raise an explicit, age-independent warning on any
change, and a website configured with a required threshold refuses bundles for any other.

\section{Security Analysis}
\label{sec:security}

Throughout, the hash functions are the Poseidon2 constructions of Section~\ref{sec:prims},
the underlying permutation is modeled as a public random permutation, SHA3/SHAKE are
modeled as random oracles \cite{bellrog}, $\negl$ denotes a negligible function of the
security parameter, $q$ bounds the adversary's oracle queries, $\tau$ is the repetition
count, and $\adv$ is a PPT (or QPT, Section~\ref{sec:pq}) adversary. The \emph{published
tree} means the complete depth-$d$ binary tree determined by the published leaf sequence,
with every unoccupied position holding the empty leaf $E$ and every unoccupied subtree
holding its cached digest; every position at every level is thus well defined.

\subsection{Unforgeability (G1, R1)}

\begin{theorem}
\label{thm:unforge}
Assume view commitments are binding (SHA3-256 as a random oracle), $\LeafHash$ and
$\Compress$ are collision and preimage resistant, challenges are derived by Fiat-Shamir in
the random oracle model, and the root the website verifies against comes from a bundle
whose ML-DSA-65 signature it checked (so that ``published tree'' is well defined; without
EUF-CMA of the signature scheme a forged bundle would let an adversary choose the root). If any coalition of citizens, websites, and
holders of the GA's public information causes an honest website to accept a token, then
except with probability at most $q(2/3)^\tau + O(q^2)\,2^{-256}$ there exists an
extractable witness $(s, \mathrm{path})$ such that $\LeafHash(s)$ is a leaf of the
published tree; and if no coalition member was enrolled with that leaf, the extraction
additionally yields a preimage or collision of the hash functions.
\end{theorem}

\begin{proof}
Fix an accepting token $(k, e, D, \pi)$, let $\mathit{root}$ be the value the website
verified against, and let $c = \CtxHash(D, e)$ be the context the website recomputed.

\emph{Step 1 (commitments bind the views).} Each repetition carries three SHA3-256
commitments over (seed, view). Producing two distinct openings for any single commitment
requires a collision in the random oracle, so over all of $\adv$'s $q$ queries this
happens with probability at most $q^2 2^{-256}$. Except in that event, at the moment a
transcript is first submitted to the challenge oracle, every repetition of that transcript
is bound to one fixed triple of views: two of them determined entirely by their seeds, the
third by its seed together with its committed input shares and multiplication outputs.

\emph{Step 2 (a false statement corrupts at least one pair per repetition).} Consider one
repetition, its bound triple of views, and the three per-party output-share vectors
included in the transcript. Party indices run over $\{0,1,2\}$ and are read modulo 3
throughout. Call the pair $(i, i{+}1)$ \emph{consistent} if the verifier's
deterministic recomputation for challenge $i$ succeeds: party $i$'s view re-derives
exactly from its seed (and committed inputs where applicable), every multiplication gate
of party $i$ is consistent with the tapes of parties $i$ and $i{+}1$ and the committed
gate outputs of party $i{+}1$, and both parties' circuit outputs equal the transcript's
output shares. If all three pairs are consistent, the three input-share vectors jointly
define a witness $w$ (the coordinate-wise sum of the shares), and by induction over the
circuit's gates every wire's three shares sum to the native evaluation of that wire on
$w$; since the verifier also checks that the output shares sum to $(k, \mathit{root},
0^d)$, $w$ satisfies $S(\mathit{root}, k, c)$. Taking the contrapositive: if the bound
views embed no witness for $S$, at least one of the three pairs is inconsistent, and the
verifier accepts the repetition only when the challenge lands on one of the at most two
consistent pairs, probability at most $2/3$.

\emph{Step 3 (Fiat-Shamir).} Challenges for a transcript $T$ are the random oracle's
output on $T$. For any single $T$ bound as in Step 1 to views lacking a witness, the
probability that its challenge vector avoids every inconsistent pair across all $\tau$
repetitions is at most $(2/3)^\tau$. A union bound over the adversary's $q$ oracle queries
bounds its total success probability by $q\,(2/3)^\tau$, which for $\tau = 219$ is below
$q\,2^{-128}$. The union is over transcripts rather than repetitions because the verifier
derives all $\tau$ challenges from the single hash of one transcript and rejects unless
every repetition passes, so the adversary gets one draw of the whole challenge vector per
query, not one draw per repetition.

\emph{Step 4 (extraction and reduction to hash security).} In the remaining case, some
accepted transcript's bound views contain a witness $(s, (b_i, \sigma_i)_{i<d})$ for $S$.
The booleanity outputs force each $b_i \in \{0,1\}$ exactly, so $\MerkleFold$ performs a
genuine root recomputation along the path selected by the bits. Compare this recomputation
to the published tree, level by level from the root down: either every node produced
equals the published node at the corresponding position, in which case $\LeafHash(s)$
equals the published leaf at index $\sum_i b_i 2^i$; or there is a highest level at which
the recomputed child pair differs from the published children while the parent agrees,
which exhibits a collision in $\Compress$. In the first case, if the indexed position
holds the empty leaf $E$, then $(s, \mathit{tag}_{\mathrm{leaf}}, 0, 0, 0)$ and
$(0,0,0,0,\mathit{tag}_{\mathrm{empty}},0,0,0)$ are distinct permutation inputs (they
differ in the tag slot, since $\mathit{tag}_{\mathrm{leaf}} \ne
\mathit{tag}_{\mathrm{empty}}$, whatever $s$ is) with equal truncated outputs, a collision
across the domain-separated instances; if the position
holds an honest citizen's commitment, $s$ is a preimage of that commitment. Each event
contradicts the assumed hash security except with negligible probability, so an accepted
token implies knowledge of an enrolled secret, within the stated bounds.
\end{proof}

\subsection{Anonymity and collusion resistance (G2, R2, R4)}

\begin{definitionx}[Anonymity game]
The adversary controls the GA and all websites, runs enrollment for citizens of its choice
plus two honest wallets $W_0, W_1$ which it enrolls (learning their leaves and indices,
but not their secrets). It then queries a token oracle at most $q$ times: a query
$(j, D, e, n)$ with $j \in \{0,1\}$ returns the token $W_j$ would produce for $(D, e, n)$
against the current published root. At challenge time $\adv$ picks $(D^*, e^*)$ at which
it has made no token query for either honest wallet, together with a nonce $n^*$ of its
choice; the challenger flips a fair coin $\beta \in \{0,1\}$, and $W_\beta$ produces a
token for $(D^*, e^*, n^*)$ which $\adv$ receives. $\adv$ outputs a guess $\beta'$. The
system is anonymous if $|\Pr[\beta' = \beta] - 1/2| \le \negl$.
\end{definitionx}

\begin{theorem}
\label{thm:anon}
In the random oracle and random permutation model, the protocol is anonymous in the above
game, with adversarial advantage at most $O(q)\,2^{-253}$.
\end{theorem}

\begin{proof}
We proceed by a hybrid argument over the challenge token.

Hybrid $H_0$ is the real game.

Hybrid $H_1$ replaces the proof $\pi$ by a simulated proof. The simulator samples the
challenge vector first; then, for each repetition, it constructs the two views to be
opened exactly as an honest prover would (uniform tapes expanded from fresh seeds, input
shares tape-derived or explicit as the party indices dictate), completes their
gate-by-gate values so that all pair checks and both output-share vectors are consistent
with the public outputs, commits honestly to those two views, commits to a uniformly
random string in place of the third view, and programs the random oracle so that the
resulting transcript maps to the pre-sampled challenges. The opened data in $H_1$ is
distributed identically to $H_0$: in both worlds the two opened views consist of uniform,
independent tape values and of shares whose joint distribution is fixed by the same
consistency constraints, because any two of the three additive shares of every wire are
uniform and independent of the witness. Two failure events separate the hybrids: the
adversary queries the random oracle on the preimage of the unopened commitment, which
contains a fresh 256-bit seed it never sees, probability at most $q\,2^{-256}$; or the
transcript to be programmed was already queried, which requires guessing one of the fresh
seeds inside the commitments, again at most $q\,2^{-256}$.

Hybrid $H_2$ replaces $k = \PRF(s_\beta, c^*)$, where $c^* = \CtxHash(D^*, e^*)$ and
$s_\beta$ is $W_\beta$'s secret, by a uniform random digest. Since every proof is
simulated after $H_1$, the adversary's entire view of $s_\beta$ is a list of digests:
$\mathit{leaf}_\beta = \Trunc\bigl(P(s_{\beta,0}, \ldots, s_{\beta,3},
\mathit{tag}_{\mathrm{leaf}}, 0, 0, 0)\bigr)$ from enrollment, the pseudonyms
$\PRF(s_\beta, c_1), \ldots, \PRF(s_\beta, c_m)$ returned by its $m \le q$ token queries,
and the challenge $k$. Write $y = P(s_{\beta,0}, \ldots, s_{\beta,3},
\mathit{tag}_{\mathrm{prf}}, 0, 0, 0)$ for the shared inner image.

Two observations give the hybrid. First, $s_\beta$ is uniform over $\Fp^4$ (min-entropy
above $255$ bits), so the leaf state and the inner PRF state are distinct points of $P$,
differing in the tag slot; with $P$ a public random permutation their images are
independent and uniform unless $\adv$ queries $P$ or $P^{-1}$ at a point whose rate half
is $s_\beta$, which each of its $q$ queries does with probability at most $2^{-255}$.
Second, the game forbids a token query at $(D^*, e^*)$, so $c_i \ne c^*$ for every $i$
except when $\CtxHash$ collides on distinct inputs, probability at most $q\,2^{-255}$;
the outer permutation calls are then evaluated at the $m+1$ distinct points $y \boxplus
c_1, \ldots, y \boxplus c_m, y \boxplus c^*$, whose images are jointly uniform and
independent unless $\adv$ queries $P$ at one of them, and $y$ is itself unqueried and
uniform. Conditioning on neither event, $k$ is uniform and independent of everything else
$\adv$ holds, exactly as in $H_2$. Hence $H_1$ and $H_2$ differ by at most
$O(q)\,2^{-253}$. Note that the $m$ prior pseudonyms are what make this step necessary:
they are the only other place $s_\beta$ is used, and they are handled rather than
excluded.

In $H_2$ the challenge token is (uniform $k$, $e^*$, $D^*$, simulated $\pi$) and is a
sample from a distribution that does not depend on $\beta$, so the advantage in $H_2$ is
zero. Summing the hybrid distances proves the theorem.
\end{proof}

\begin{corollary}[What a full-collusion coalition holds]
The coalition of the GA and all websites possesses: the registry (citizen $\to$ leaf), the
public tree, and per site a set of pseudonyms $k$ with session metadata. Theorem~\ref{thm:anon}
compares two wallets; a standard hybrid over the enrolled population, replacing one
citizen's pseudonyms by uniform digests at a time, extends it to the statement that the
joint distribution of all pseudonyms is computationally independent of the registry, at a
cost linear in the number of enrolled citizens. There is additionally no traffic to correlate: the protocol defines no
message from citizen or website to the GA after enrollment, so timing and intersection
attacks in the style of mix-network analysis have no government-side observation to work
with. The residual signal is the website-side one that any web service has (IP address,
timing), which is out of scope by Section~\ref{sec:model} and unchanged by this system.
\end{corollary}

\subsection{Cross-site and cross-epoch unlinkability (G3, R3)}

\begin{theorem}
\label{thm:unlink}
Let $(D, e) \ne (D', e')$ be distinct domain-epoch pairs and consider the following game.
The adversary controls the GA and all websites and enrolls two honest wallets with secrets
$s_0, s_1 \in \Fp^4$; the challenger flips $\beta \in \{0,1\}$ and returns the token of
$s_0$ at $(D, e)$ together with the token of $s_\beta$ at $(D', e')$, that is, the
pseudonyms $k = \PRF(s_0, \CtxHash(D, e))$ and $k' = \PRF(s_\beta, \CtxHash(D', e'))$ with
their proofs. Under the model of Theorem~\ref{thm:anon}, no adversary decides whether the
two tokens come from one secret ($\beta = 0$) or two ($\beta = 1$) with advantage above
$O(q)\,2^{-253}$, where $q$ bounds its oracle and permutation queries.
\end{theorem}

\begin{proof}
Write $c = \CtxHash(D, e)$ and $c' = \CtxHash(D', e')$. As $\CtxHash$ is SHAKE256 with
rejection sampling, modeled as a random oracle, its outputs on the two distinct inputs
$(D, e) \ne (D', e')$ are independent and uniform on $\Fp^4$, so $c \ne c'$ except with
probability $p^{-4} < 2^{-255}$. This is a collision between two fixed inputs, so the
bound does not depend on how many queries the adversary makes. Apply the hybrid sequence of
Theorem~\ref{thm:anon} to both tokens simultaneously: the simulation step is per-token and
unchanged; the replacement step substitutes both pseudonyms by uniform values, justified
by the same random permutation argument, since the two outer PRF calls evaluate $P$ at
distinct points (their rate halves differ by $c$ versus $c'$) and are therefore jointly
uniform and independent unless the adversary queries $P$ at a point whose rate half is
$s_0$ or $s_1$, which each of its $q$ queries does with probability at most $2^{-255}$.
When $\beta = 0$ the two outer calls share one inner image and when $\beta = 1$ they use
two, but in both cases the evaluation points are distinct, which is why the case
distinction never surfaces in the adversary's view. In the
final hybrid both tokens are independent uniform samples, a distribution identical under
$\beta = 0$ and $\beta = 1$, so the advantage there is zero; summing the hybrid distances
gives the stated bound. Within a single $(D, e)$, linkability of two visits is intentional
($k$ is deterministic) and disclosed as the pseudonymity property.
\end{proof}

\subsection{Replay resistance (G4)}

\begin{theorem}
\label{thm:replay}
Let $(k, e, D, \pi)$ be a token produced by an honest wallet and accepted by site $D$ with
nonce $n$. Except with probability at most $(q^2 + 3\tau)\,2^{-256} + 3^{-\tau}$, the
token verifies at no other tuple $(D', e', \mathit{root}', k', n') \ne (D, e,
\mathit{root}, k, n)$, and at $(D, n)$ itself at most once.
\end{theorem}

\begin{proof}
Write the proof as a transcript body $T$ (all commitments and output shares) plus openings
$O$ for the challenge vector $\chall \in \{0,1,2\}^\tau$, derived as $\chall = \RO(\mathrm{prefix}
\,\|\, \mathit{ctx} \,\|\, T)$, where $\RO$ is the random oracle instantiating the
Fiat-Shamir hash (SHA3-256), $\mathrm{prefix}$ is the protocol's fixed domain-separation
string, and $\mathit{ctx} = \Enc(D, e, \mathit{root}, k, n)$; $\chall_i$ denotes the $i$-th
entry of $\chall$ and party indices are taken modulo 3. Honest generation draws the three
per-repetition seeds independently and uniformly, so within a repetition the three seeds
are pairwise distinct except with probability at most $3\,2^{-256}$, hence at most
$3\tau\,2^{-256}$ over all $\tau$ repetitions. A verifier for
a different tuple derives $\chall' = \RO(\mathrm{prefix} \,\|\, \mathit{ctx}' \,\|\, T)$. The two contexts encode
distinct tuples, so $\mathit{ctx}' \ne \mathit{ctx}$ as bitstrings by
Lemma~\ref{lem:enc}, and $\chall' = \chall$ only if the random oracle
collides on distinct inputs, probability at most $q^2 2^{-256}$ over the adversary's
queries; otherwise $\chall'$ is uniform on the $\tau$-fold challenge space and independent of
$\chall$. The openings $O$ contain, for each repetition $i$, the seeds of parties $\chall_i$ and
$\chall_i + 1$ and the multiplication outputs of party $\chall_i + 1$. If $\chall'_i \ne \chall_i$ at any
repetition, the verifier interprets the provided seeds as belonging to parties $\chall'_i$ and
$\chall'_i + 1$ and checks them against the commitments at those indices; each commitment
contains its own party's fresh 256-bit seed, and those seeds are pairwise distinct, so a
seed committed at one index opens the
commitment at a different index only via a random-oracle collision, already accounted for.
Acceptance therefore requires $\chall' = \chall$ exactly, an event of probability $3^{-\tau}$, which
for $\tau = 219$ is below $2^{-347}$. Re-presentation at $(D, n)$ itself is blocked by the
verifier's local consumed-nonce state, which needs no coordination with any other party;
and a fresh nonce at the same site changes $\mathit{ctx}$, reducing to the case above.
\end{proof}

\paragraph{Remark.} A wallet that deviates from honest randomness (for example, giving all
three parties identical seeds) can craft tokens of its own that verify under several
challenge vectors and are therefore replayable by third parties. This weakens nothing but
the misbehaving wallet's own binding: it grants adulthood attestations only with the
wallet owner's active cooperation, which is exactly the voluntary-sharing case of
Section~\ref{sec:limits}.

\subsection{Threshold isolation (multi-threshold deployments)}
\label{sec:isolation}

This subsection proves the claim of Section~\ref{sec:thresholds}, mechanism 1: a token is
usable against exactly the threshold it was produced for, and testing it against any other
threshold produces a verdict that carries no information about the visitor.

\begin{proposition}[Threshold isolation]
\label{prop:isolation}
Let $a \ne b$ be age thresholds with published roots $\mathit{root}_a$, $\mathit{root}_b$,
and let $(k, e, D, \pi)$ be a token produced by an honest wallet for threshold $a$ at
domain $D$ and epoch $e$ with website nonce $n$, that is, a proof of $S(\mathit{root}_a,
k, c_a)$ with $c_a = \CtxHash(D, e, a)$, transcript-bound as in
Section~\ref{sec:thresholds}. For every nonce $n'$, a verifier that checks this token
against threshold $b$, that is, against $(\mathit{root}_b, c_b = \CtxHash(D, e, b), n')$,
accepts with probability at most $(q^2 + 3\tau)\,2^{-256} + 3^{-\tau}$, where $q$ bounds
the adversary's random-oracle queries and $\tau$ is the repetition count
(Section~\ref{sec:security}). The bound holds uniformly over all wallets; in particular
it does not depend on whether the wallet's leaf lies in $\tree_b$.
\end{proposition}

\begin{proof}
The threshold's encoding into the Fiat-Shamir context is injective. Both sides derive the
context from their own $(D, a)$ by forming the string $D \,\|\, \texttt{\#} \,\|\, a$ and
using it in place of the bare domain, and \texttt{\#} occurs in no hostname, so the map
$(D, a) \mapsto D \,\|\, \texttt{\#} \,\|\, a$ is injective and $a \ne b$ forces distinct
domain fields. Hence the prover's context bitstring $\mathit{ctx}_a$, encoding $(D
\,\|\, \texttt{\#} \,\|\, a,\, e,\, \mathit{root}_a,\, k,\, n)$, and the $b$ verifier's
$\mathit{ctx}_b$, encoding $(D \,\|\, \texttt{\#} \,\|\, b,\, e,\, \mathit{root}_b,\, k,\,
n')$, differ as bitstrings for every choice of $n'$: they already differ in the domain
field, whether or not $\mathit{root}_a$ and $\mathit{root}_b$ happen to coincide.
Lemma~\ref{lem:enc} then gives $\mathit{ctx}_a \ne \mathit{ctx}_b$, and
Theorem~\ref{thm:replay} applies verbatim with the target tuple taken to be the $b$
verifier's tuple, yielding the stated bound. Its derivation uses only the honest
generation of $\pi$ (independent uniform seeds) and random-oracle properties; at no point
does it reference the wallet's witness or the contents of $\tree_b$. Hence every wallet's
token is rejected by the $b$ verifier with probability at least $1 - (q^2 +
3\tau)\,2^{-256} - 3^{-\tau}$: the verdict has the same distribution whether or not the
holder is enrolled in $\tree_b$, and a test whose output distribution does not depend on
the holder conveys nothing about the holder. That is the claimed uniformity.
\end{proof}

\subsection{Post-quantum security (G5, R5)}
\label{sec:pq}

Every security property above reduces to: preimage resistance and PRF security of
Poseidon2-based functions with 256-bit capacity and digests, which Grover \cite{grover}
reduces to 128 bits; collision resistance of the same functions, which deserves a separate
accounting given below; random-oracle behavior of SHA3/SHAKE, where soundness of the
Fiat-Shamir transformation against quantum adversaries follows from the
measure-and-reprogram analyses for commit-and-open protocols in the quantum random oracle
model \cite{qromfs}, whose polynomial query losses are absorbed by our $2^{-128}$
classical soundness margin; and EUF-CMA of ML-DSA-65 under Module-LWE and Module-SIS
\cite{regev,dilithium}, NIST category 3. No discrete-log, factoring, or pairing assumption
appears anywhere. Note the contrast with pre-quantum anonymous-credential deployments:
here even the issuer's signature protects only root authenticity; a future break of ML-DSA
would permit root forgery (a soundness attack, mitigated by re-signing with a successor
scheme such as the hash-based SPHINCS+ family \cite{sphincs}) but could never
retroactively deanonymize anyone, because no signature ever touches a per-citizen value.

\paragraph{Collision resistance under quantum attack.} Grover is the wrong tool for
collisions and we do not claim 128-bit quantum collision resistance. Our 256-bit digests
give 128-bit classical collision resistance; the BHT algorithm \cite{bht} finds collisions
in about $2^{85}$ quantum operations, but needs comparable quantum-accessible memory, and
under any cost metric that charges for that memory it does not beat classical parallel
search \cite{bernsteincoll}. Two observations bound the consequences. First, no privacy
property rests on collision resistance of $\LeafHash$ or $\Compress$: Theorems
\ref{thm:anon} and \ref{thm:unlink} use only the random-permutation behavior of $P$ and
the min-entropy of $s$, so even a full collision break costs no anonymity, and it is
anonymity that cannot be repaired after the fact. Second, the properties that do rest on
collision resistance are repairable in place, and it is worth separating which function
each one needs. Unforgeability (Theorem~\ref{thm:unforge}) needs collision resistance of
$\LeafHash$ and $\Compress$, both Poseidon2 with 256-bit digests. The separation of
contexts in Theorem~\ref{thm:unlink} needs $\CtxHash$ to be collision resistant, which is
a property of SHAKE256 truncated to $\Fp^4$, not of Poseidon2; the two are widened
together, since both outputs are four field elements. Neither degrades below the $2^{85}$
figure, and that figure holds only under the optimistic memory model. A jurisdiction that
judges the margin insufficient widens both digests to six field elements, about 384 bits
(state width $t = 12$, the Poseidon2 parameter set whose known-answer vector we already
validate against), at a cost linear in the circuit size and with no change to the protocol
or to any proof above. Note that the injectivity of the \emph{context encoding}
(Lemma~\ref{lem:enc}), on which Theorem~\ref{thm:replay} and
Proposition~\ref{prop:isolation} rest, is a syntactic property of a byte layout and is not
affected by any cryptanalytic advance at all.

\subsection{What the system does not solve}
\label{sec:limits}

\paragraph{Voluntary lending.} An adult may run the wallet for a minor, or hand over the
device. No credential system prevents this without continuous biometric presence-testing,
which would destroy the privacy goal. Mitigations are friction and consequence: epoch
rotation limits the blast radius of a shared secret; revocation punishes detected
commercial lending (a lent secret used at scale is statistically visible to sites as one
pseudonym with implausible behavior, reportable for revocation without identifying
anyone); optional device binding (secret sealed in a secure element) raises extraction
cost. We consciously reject the addition of a per-use government authorization step: it
would add a metadata channel while, being blind, it could not actually bind uses to
persons, so it would purchase surveillance risk and buy no sharing resistance.

\paragraph{Sybil and soundness attacks by the GA.} A malicious GA can enroll fictitious
adults or refuse real ones. This is an attack on soundness and fairness, not on privacy,
and it is auditable: the tree is public and append-only, its growth must match published
demographic statistics, and each insertion can be required to carry a transparency-log
inclusion receipt.

\paragraph{Split-view roots.} A GA serving a victim a personalized root is detected by
bundle gossip (Section~\ref{sec:rootpub}); a wallet that verifies its bundle against two
independent mirrors reduces the attack to requiring collusion of the mirrors.

\paragraph{Enrollment coercion and the enrollment event itself.} The GA knows who
enrolled. Enrollment must therefore be normalized (bundled with routine eID issuance for
all adults, not requested on demand) so that holding a credential carries zero information
about intent to use it.

\section{Implementation}
\label{sec:impl}

The reference implementation is about 1{,}150 lines of dependency-light Python in
\texttt{code/python/} (\texttt{core.py}: field, Poseidon2 native and circuit evaluation,
sparse Merkle tree, prover and verifier, binary wire codec; \texttt{protocol.py}: GA,
wallet with encrypted storage, website verifier), plus a test suite and a scenario demo.
Everything in this paper, including its sources, is at
\url{https://github.com/Kidev/ZeroKnowledgeAgeVerification}. Design points of engineering
interest:

\paragraph{Engine-generic circuit.} The circuit is written once against an abstract engine
interface (const, add, sub, add-constant, mul-constant, mul) and executed by three
engines: a native engine on field elements (used by GA, wallet, and website for tree
building and expected-output computation), a prover engine on 3-tuples of additive shares
that appends every multiplication-gate output share to per-party views, and a pair engine
that replays the two challenged parties for the verifier, recomputing the first and
consuming the second's committed outputs. The native permutation is generic in the state
width, so the Poseidon2 designers' published $t = 12$ Goldilocks known-answer vector is
checked through that same routine rather than through a second implementation written for
the test; the circuit engine's permutation is then cross-checked against it at $t = 8$ on
random states, and the $t = 8$ constants are extracted verbatim from the designers'
reference implementation. The chain from published vector to proving circuit therefore has
no link that a test could satisfy while production code drifted.

\paragraph{Determinism and view minimization.} Party tapes are expanded from 32-byte seeds
by SHAKE256 in counter mode with rejection sampling (exactly uniform field elements);
parties 0 and 1 take their input shares from their tape prefixes, so an opening transmits
only two seeds, the neighbor's multiplication outputs, and (when party 2 is opened) its
input shares. Commitments are SHA3-256 over seed and view; the verifier recomputes the
challenged party's view rather than receiving it, so a corrupted committed view is caught
with certainty, not probability.

\paragraph{Wire format and hardening.} Proofs serialize to a versioned binary format with
fixed layouts derived from the depth and repetition count, with every multi-byte integer
little-endian independently of host byte order, so the format is a specification rather
than an artifact of one architecture. The verifier is total: it validates commitment
lengths, share ranges, view lengths, and exact tape consumption \emph{before} any value
reaches a hash input, returns false on every malformed input rather than raising, and
treats the transcript hash as the single source of challenge truth. The wire decoder
bounds the depth and repetition count declared in the header, and cross-checks them
against the blob length, before sizing any allocation from them; a decoder that trusts an
attacker-supplied length is a memory exhaustion primitive. Both properties are covered by
the test suite.

\paragraph{Failure containment.} The wallet secret is 4 uniform field elements generated
by the OS CSPRNG, stored only under scrypt-derived AES-256-GCM, and used solely inside
$\LeafHash$ and $\PRF$; there is no code path that serializes it into any protocol
message. Domain strings are rejected if they contain the context separator, which is what
discharges the hypothesis of Lemma~\ref{lem:enc} rather than assuming it.

\paragraph{Threshold binding.} In multi-threshold deployments
(Section~\ref{sec:thresholds}) the threshold enters the signed root bundle, the proof
context (encoded as the context string \texttt{domain\#threshold}; \texttt{\#} cannot
occur in a hostname, so single-set deployments remain byte-compatible), and the wallet's
per-domain pin store, persisted inside the encrypted wallet file. The wallet exposes the
requested threshold for display before proving; a threshold change raises a dedicated
warning type that the UI must surface, and proving proceeds only under an explicit
override flag. The warning check runs before any membership-dependent computation, so
wallet behavior on a threshold change is independent of the holder's age. If the holder
turns out to hold no leaf in the requested accumulator, the wallet declines locally
instead of emitting a proof that cannot verify: a failing proof would announce ``this
holder is below that threshold'', whereas a local decline is indistinguishable, from the
site's side, from a user who simply refused. The site sees two outcomes, token or no
token, rather than three.

\paragraph{Native port.} A Rust implementation of the field, permutation, tree, circuit
engines, prover, and verifier (\texttt{code/rust/}, about 1{,}250 lines, sharing the
Fiat-Shamir transcript and binary wire format) is pinned to the reference by known-answer
tests extracted from the Python code and by wire-format interoperability: each
implementation verifies, at full parameters, proofs produced by the other, so the two
codebases cannot drift apart silently. Its measured performance is reported in
Section~\ref{sec:perf}.

\section{Performance}
\label{sec:perf}

\subsection{Measured}

Single core of an AMD Ryzen~9 7950X, CPython 3.14, pure Python (no native
extensions), $\tau$ repetitions, depth-$d$ tree:

\begin{center}
\begin{tabular}{rrrlrrr}
\toprule
depth $d$ & capacity & $\tau$ & soundness & prove & verify & proof size \\
\midrule
20 & 1.05\,M  & 219 & $< 2^{-128}$ & 3.1\,s & 2.0\,s & 14.3\,MB \\
27 & 134\,M   & 219 & $< 2^{-128}$ & 4.1\,s & 2.6\,s & 18.7\,MB \\
32 & 4.29\,B  & 137 & $< 2^{-80}$  & 3.0\,s & 1.9\,s & 13.6\,MB \\
32 & 4.29\,B  & 219 & $< 2^{-128}$ & 4.7\,s & 3.0\,s & 21.8\,MB \\
\bottomrule
\end{tabular}
\end{center}

Enrollment is one hash and one $O(d)$ tree insertion (microseconds to milliseconds); root
bundle signing is one ML-DSA-65 signature (3{,}309 bytes); wallet encryption and
decryption are dominated by scrypt at about 100\,ms.

The native Rust port of the same backend (Section~\ref{sec:impl}), producing
byte-identical proofs, measures on the same machine, on which CPython takes 4.7\,s to
prove and 3.0\,s to verify at $d = 32$, $\tau = 219$:

\begin{center}
\begin{tabular}{rrrlrrr}
\toprule
depth $d$ & capacity & $\tau$ & soundness & prove & verify & proof size \\
\midrule
32 & 4.29\,B  & 137 & $< 2^{-80}$  & 0.20\,s & 0.12\,s & 13.6\,MB \\
32 & 4.29\,B  & 219 & $< 2^{-128}$ & 0.33\,s & 0.18\,s & 21.8\,MB \\
\bottomrule
\end{tabular}
\end{center}

That is roughly fifteen times faster than CPython on the same machine, from a
straightforward port with no low-level optimization, and already under a second at full
parameters on one core; repetitions additionally parallelize across cores.

\subsection{Cost model}

The circuit has $344(d+3) + 5d$ multiplication gates; prover work is $3\tau$ times that,
verifier work $2\tau$ times, and proof size is dominated by the transmitted neighbor views
at 8 bytes per gate per repetition. All costs are linear in both $d$ and $\tau$ and
embarrassingly parallel across repetitions.

\subsection{Deployment path}
\label{sec:deploypath}

The measured numbers above show the reference backend already crossing the sub-second bar
in its native port. Two further steps, neither changing the protocol or the statement,
close the remaining gap to interactive-latency deployment. (i) Native optimization of the
included Rust prover: vectorized Keccak for tape expansion, in-place circuit evaluation,
precomputed share layouts, and parallelism across repetitions bring proving to the low
hundreds of milliseconds and below on commodity hardware. (ii) Replacing the backend
with a hash-based STARK or a modern MPC-in-the-head system (KKW/Picnic3-style
preprocessing, or VOLE-in-the-head) over the identical Poseidon2 circuit reduces proofs to
the 100 to 250\,KB range and verification to tens of milliseconds, with the same
no-trusted-setup, symmetric-only trust profile. Because verification is website-local,
aggregate system capacity is simply the sum of website capacity; the government's serving
burden is a static file (the bundle and tree deltas) on a CDN. There is no
per-verification government cost at any scale, which is not an optimization but the point.

\section{Deployment Considerations}
\label{sec:deploy}

\paragraph{Rollout.} Enroll by cohort alongside routine national eID operations so
possession is universal among adults and signals nothing. Publish the tree and bundles
from day one with independent mirrors and an append-only audit log. Ship the wallet as
reproducibly built open source in app stores and as a browser extension; the proving key
material is nil (no trusted setup), so distribution integrity is the only supply-chain
concern.

\paragraph{Website integration.} A site embeds the bundle-verification key, fetches
bundles on a schedule, issues nonces, and calls a verifier library; the entire integration
is stateless except nonce and pseudonym storage. Regulators can certify libraries rather
than audit data flows, because there are no data flows.

\paragraph{Legal framework.} Legislation should (i) authorize commitment-based enrollment,
(ii) recognize a verified token as compliance with age-gating duties, discharging website
liability, (iii) prohibit websites from demanding any additional identity signal when a
token is presented, including proofs against any age threshold above the one their duty
requires (Section~\ref{sec:thresholds}), and (iv) mandate the transparency obligations of
Section~\ref{sec:rootpub}. Point (iii) is a liability backstop rather than a load-bearing
assumption: the threshold mechanism of Section~\ref{sec:thresholds} already makes escalation
self-defeating, so the guarantee does not depend on the prohibition.

\paragraph{Governance of the one central artifact.} The root is the sole trusted object.
Its signing key should live in an HSM with a published rotation and successor-algorithm
policy (Section~\ref{sec:pq}), and bundle history should be anchored in at least one
system outside government control.

\section{Layman's Explanation}

Suppose the government publishes a long public list of scrambled codes, one per adult.
The
scrambling is one way: only your wallet, an open-source app running on your own device,
knows the secret that produces your code, a secret it generated by itself, offline,
without contacting any server. The government knows it placed your code on the list; that is
the entirety of what it will ever know.

When a website requires proof of age, the wallet produces a zero-knowledge proof, a
cryptographic technique established in the 1980s and in production use today. Think of
the government's list as a single closed chain made of millions of padlocks, one per
adult,
and imagine you are wearing a mask. You hold the key to exactly one of those padlocks. To
prove you are an adult, you do not point at your padlock and you do not show your face:
you simply open the chain. A chain can only be opened by unlocking one of its padlocks,
so everyone watching is convinced you hold a genuine key; the mask (here, the
zero-knowledge layer) guarantees that nobody can tell which padlock you unlocked, and
therefore who you are. The demonstration is repeated a few hundred times with fresh
randomness, leaving someone with no key less than one chance in $2^{128}$ of faking it.
The website learns one fact, an adult is present, and receives a pseudonym valid only on
that site and for one month, enough to recognize a returning visitor; pseudonyms on
different sites are mathematically unrelated.

For the reader who assumes the worst of the government: after enrollment, your wallet
never communicates with it again, directly or indirectly. When you visit a site, no
message exists that could be intercepted, logged, demanded, or leaked, because none is
sent. Nobody needs to be trusted to delete records of your visits; those records cannot
come into existence. The mathematics is public, the software is open to inspection and
reproducible builds, and the government's single publication is a list anyone can audit
for dishonest edits. A fully hostile government could stop adding adults or strike people
from the list, but only in public, where every change is visible. What it cannot do, with
any computer including a quantum one, is learn that you visited a particular site. The
website, for its part, knows only that some adult presented a valid proof.

\section{Conclusion}

The debate around online age verification presents privacy and child protection as
opposed. They are opposed only under bad architectures. By reducing the government's role
to publishing one signed list of scrambled commitments and pushing all verification into
local mathematics between citizen and website, this design achieves the strong form of
both properties simultaneously: websites obtain a sound adulthood attestation, and the
mapping from visits to citizens does not exist as data anywhere, under collusion, under
compulsion, and under quantum attack. The construction uses only standardized or
academically established components: Poseidon2, MPC-in-the-head zero knowledge, SHA-3,
and ML-DSA. The reference implementation accompanying this paper demonstrates every
claimed behavior and every analyzed attack. What remains is engineering (faster proof
backends, wallet distribution) and legislation (mandating that a token is enough). Neither
remaining task requires trusting anyone with the one thing this system was built to make
untrustable.

\bibliographystyle{plain}
\bibliography{sources}

\end{document}
